\section{Limitations}

\paragraph{Model and trigger coverage.}
The study covers multiple model families, but it remains centered on one LoRA poisoning setup and one trigger family.
Future work should vary trigger syntax, poison ratio, target behavior, adapter placement, and fine-tuning recipe.
The Llama result already suggests that model architecture can substantially change the attainable safety-utility frontier.

\paragraph{Evaluation protocol.}
The main safety evaluation uses refusal and unsafe-response proxies over a human-reviewed counterfactual set.
This protocol is useful for controlled large-scale comparison, but stronger judges or human audits remain important for borderline generations, partial compliance, and ambiguous refusals.
The Llama MMLU values are reported from a separate legacy-prompt MMLU protocol, which should be kept explicit when comparing utility across model families.

\paragraph{Incomplete operator space.}
\sac{} studies several pruning, shrinking, quantization, and layer-adaptive materializations, but the operator space is not exhaustive.
Additional operators, joint optimization of gates and materialization, or calibration with stronger safety judges may improve the frontier, especially for Llama-like boundary cases.

\paragraph{Incomplete defense baselines.}
The current baselines mainly establish that the effect is not caused by compression pressure alone.
They do not establish superiority to the strongest existing backdoor-removal defenses.
Completed clean-adapter, no-adapter, merged-adapter, representation-editing, refusal-tuning, and adapter-specific backdoor-removal baselines are still required before making a comparative-defense claim.
These runs are queued in the supplement harness but are not yet draft results.

\paragraph{Over-refusal.}
The Qwen4B results demonstrate strong attack suppression but substantial triggered-benign over-refusal.
This is a real deployment cost, not a cosmetic artifact.
The method should therefore be viewed as a frontier-improving framework rather than a guaranteed no-cost defense.

\paragraph{No universal recipe.}
The results do not support a universal compression setting for all adapters and base models.
They instead show that security-aware selection can be highly effective in some regimes and only partially effective in others.
Understanding which architectural or adapter properties predict this frontier remains an open problem.
